% 第 47 章。数字与教学模型：checks/check_ch47.py；图：figsrc/fig_ch47.py。
% 史料：oil_facts（取数脚本 data/fetch/oil47.py）；数据：oil_pinksheet、oil_wdi、oil_pwt（同一脚本）。
\chapter{石油危机与滞胀}\label{ch:oil}

\begin{chapterguide}
\guidequestion{油价冲击为什么同时推高通胀和失业？产油国为什么会被“资源诅咒”？}
\begin{guidegoals}
\item 用总供给冲击分析 1973 年与 1979 年的油价冲击；
\item 比较各国的政策应对（迁就或不迁就）与结果；
\item 解释石油美元回流怎样埋下拉美债务危机；
\item 推导荷兰病模型（Corden \& Neary 1982），并用产油国的案例说明资源诅咒的争论。
\end{guidegoals}
% 据: bj_1016, b80_real, b80_painful, b80_unanimous, b75_surplus, b80_burdened, cn_cover, cn_date, sw_abstract
\guideprereq{\chref{ch:naturalrate}\enspace\nameref{ch:naturalrate}；\chref{ch:floating}\enspace\nameref{ch:floating}}
\end{chapterguide}

1973 年 10 月 18 日中午刚过 12 点，国务卿基辛格（Henry Kissinger，1923—2023）在国务院主持例会\cite{frus47,loc47}。12 天前，埃及与叙利亚突然进攻以色列，战争开始了\cite{frus47}。会上谈到石油。主管近东事务的助理国务卿说，油价从每桶 3 美元涨到 5 美元出头，石油公司只能接受\cite{frus47}。桶是计量原油的体积单位。
% 据: f220_time, f209_war, f220_sisco, f220_accept, fp_sisco, loc_kissinger

基辛格追问，明年涨到 10 美元，也得接受吗？他接着问，涨到 20 美元呢，“总有该抵抗的时候吧”\cite{frus47}。常务副国务卿答，到了一定价位，煤之类的能源会比石油便宜\cite{frus47}。基辛格说，局面只会更坏，他的直觉是对抗宜早不宜迟\cite{frus47}。主管经济事务的副国务卿反问，对抗是什么意思？政府里还没有人想出应对这些要求的办法\cite{frus47}。
% 据: f220_ten, f220_twenty, f220_coal, fp_rush, f220_early, f220_casey, fp_casey

第二天，中央情报局（CIA）的经济研究处写了一份备忘录\cite{frus47}。它写道，阿拉伯产油国每月至少减产 5\%，美国每天只从它们那里得到约 160 万桶石油\cite{frus47}。西欧的石油消费有 70\% 靠阿拉伯国家，日本的供给有 40\% 靠它们\cite{frus47}。它估计美国与日本的石油进口账单各要增加约 30 亿美元，西欧接近 80 亿美元\cite{frus47}。进口国的贸易收支会急剧恶化，CIA 却认为对国内经济的影响“多半很小”\cite{frus47}。它的理由是，进口的石油不与本国的生产竞争，起初不会有相应的就业损失\cite{frus47}。
% 据: f223_date, f223_cut, f223_us, f223_dep, f223_bills, f223_slight, f223_jobs

两个多月后，判断变了。12 月 23 日，海湾产油国在德黑兰宣布，把原油的公示价格再提高一倍多，到每桶 11.65 美元，1974 年 1 月起实行\cite{boj1986b,frus47}。1974 年 1 月，CIA 估计 1974 年自由世界的石油账单要多出 700 亿美元\cite{frus47}。它说，进口国面临购买力的大量流失，在西欧与日本约合国民生产总值的 3\%\cite{frus47}。不采取扩张的措施，各国的增长都会大大放慢，产出可能下降，失业会增加。可是通胀已经很高，各国政府扩张需求时会很谨慎\cite{frus47}。油价冲击为什么同时推高通胀和失业？产油国为什么会被“资源诅咒”？
% 据: bj_1223, f273_tehran, f277_date, f277_70, f277_loss, f277_unless, f277_cautious

\section{当时的世界}\label{sec:oil-world}

禁运早在预料之中。10 月 6 日开战当天上午 9 点，华盛顿特别行动小组开会，国防部长问到石油禁运的可能，主管近东事务的副助理国务卿说可能性“非常高”\cite{frus47}。与会者一致认为，沙特阿拉伯是石油问题的关键\cite{frus47}。常务副国务卿说，美国没有应付禁运的计划，真有禁运，美国会很狼狈，欧洲人与日本人也一样\cite{frus47}。第二天傍晚，基辛格对小组说，阿拉伯人得懂得界限在哪里，不然会一点一点把美国咬死，可眼下偏偏不是教训他们的时候\cite{frus47}。
% 据: f209_wsag, f209_high, fp_atherton, f209_key, f209_noplans, f209_fix, f209_oct7, f209_nibble

谁作的决定？10 月 16 日，海湾的 6 个产油国突然宣布提高阿拉伯轻质原油的公示价格\cite{boj1986b}。每桶从 3.01 美元提到 5.12 美元，一下提高 70\%\cite{boj1986b}。公示价格是产油国定下的标价，它们从每桶油得到的收入按它计算（本书的解读）。第二天，阿拉伯石油输出国组织（OAPEC）开紧急会议，决定以 9 月的产量为基准，从 10 月起每月减产 5\%，直到以色列撤出占领地\cite{boj1986b}。日本银行说，1972 年底 OAPEC 国家的石油产量约占全世界的 32\%，出口占 55\%\cite{boj1986b}。
% 据: bj_1016, bj_1017, bj_32

沙特的立场 FRUS 有记载。石油大臣亚马尼（Ahmad Zaki Yamani，1930—2021）对法国人说，美国若不促成公正的解决，沙特会在一年里每月减产 5\%\cite{frus47,loc47}。10 月 14 日，基辛格写信给费萨尔国王（Faisal，1906—1975），说美国向以色列空运物资，并不是要与阿拉伯为敌\cite{frus47,loc47}。16 日，国王回信说，得知美国的决定他很“痛心”\cite{frus47}。从 10 月 20 日起，阿拉伯产油国相继对支持以色列的美国与荷兰全面停止出口石油\cite{boj1986b}。11 月 4 日，OAPEC 宣布 11 月的产量比 9 月减少 25\%，12 月再减 5\%\cite{boj1986b}。
% 据: f216_yamani, fp_yamani, gnd_yamani, gnd_yamani_born, f216_airlift, loc_faisal, f216_pained, bj_embargo, bj_1104

减产的规模不小。按 Hamilton 的计算，11 月阿拉伯产油国的产量比 9 月每天少了 440 万桶，相当于全世界产量的 7.5\%\cite{hamilton2011}。伊朗在禁运期间反倒增加了产量，可只抵得上一小部分\cite{hamilton2011}。国际清算银行（BIS）说，12 月底提价两天以后，阿拉伯产油国决定把产量恢复到 9 月的 85\%\cite{bis1974}。它说，这时人们看清了，主要问题在价格，不在数量\cite{bis1974}。禁运的对象只是美国与荷兰，可石油在世界市场上流通，减产与提价才落到所有进口国的头上（本书的解读）。
% 据: hm_cut, hm_offset, hm_defiance, b74_85, b74_price

德黑兰的提价是谁定的？美国驻伊朗大使在电报里说，伊朗国王巴列维（Mohammad Reza Pahlavi，1919—1980）这一年想法大变\cite{frus47,loc47}。春天他还主张压低油价，夏天美国的商品价格翻番，伊朗的通胀让他头疼\cite{frus47}。国王说，他卖油所得没法比别国少\cite{frus47}。FRUS 的编者转述，沙特在德黑兰会上“强烈反对”提价，其余的 OPEC 成员在国王带领下要每桶 10 到 14 美元\cite{frus47}。费萨尔国王认为提价会伤害世界，为此把亚马尼训了一顿\cite{frus47}。国王巴列维当天对记者说，他同情发展中国家，伊朗多出的收入有一部分会投到发达国家与发展中国家\cite{frus47}。
% 据: fp_helms, f271_changed, f271_staggered, f271_shah, f271_opposed, f271_faisal, f271_press, loc_shah

两天以后，阿拉伯产油国又宣布，1974 年 1 月增产 10\%，足量供应英国、法国、日本、西班牙等友好国家，对美国的禁运照旧\cite{frus47}。美国驻沙特大使 Akins 抱怨，德黑兰是一场灾难，科威特的 OAPEC 会议是一种侮辱\cite{frus47}。12 月 28 日，基辛格写信给沙特的外交国务大臣，称德黑兰的提价“剧烈而毫无道理”\cite{frus47}。他说，单单对美国禁运，使尼克松总统陷入不可能的处境\cite{frus47}。
% 据: f273_more, f273_dec25, f271_akins, f271_kuwait, fp_akins, f273_date, f273_tehran, fp_saqqaf, f273_nixon

美元的贬值也在背景里。日本银行说，提价的背景之一，是美元贬值等使原油出口收入缩水的倾向加强了\cite{boj1986b}。\chref{ch:floating}的结尾讲过，以美元收款的产油国，美元贬值就等于油卖得便宜了，这里有了史料的印证（本书的解读）。
% 据: bj_dollar

两次提价改变了钱的分配。BIS 说，产油国从石油公司持有股份的那部分产量中，每桶所得从 1.75 美元升到 7 美元\cite{bis1974}。同一种油运到欧洲的到岸价，从 3.65 美元升到 8.9 美元\cite{bis1974}。到岸价包括运费与保险（本书的解读）。产油国的所得翻了两番，到岸价只涨了一倍多，石油公司与运费分到的部分变小了（本书的计算）。
% 据: b74_take, b74_cif

各国受的冲击不一样。按 BIS 的统计，石油进口占能源总供给的比例，日本是 70\%，西欧平均约 60\%，美国只有 10\%\cite{bis1974}。它说，提价对消费物价的直接影响，在多数欧洲国家是 2\% 到 3\%，在日本接近 4\%\cite{bis1974}。美国与加拿大本国的能源多，价格还能压住一些，影响小一些\cite{bis1974}。
% 据: b74_dep, b74_dep_us, b74_cpi, b74_cpi_us

\figref{fig:oil_price}画出世界银行的原油年均价，是布伦特、迪拜与西得克萨斯三种原油的平均\cite{pinksheet2026}。名义价格 1973 年是每桶 2.8 美元，1974 年是 11.0 美元\cite{pinksheet2026}。实际价格按 2010 年的美元计，从 9.1 美元升到 29.2 美元\cite{pinksheet2026}。世界银行用制造品出口的单位价值指数（MUV）折算实际价格，这个指数记的是 15 个国家出口的制造品以美元计的单价\cite{pinksheet2026}。按实际价格，第一次冲击使油价一年里涨到原来的 3 倍多，此后几年慢慢回落，1978 年是 24.3 美元，第二次又从这里涨到 1980 年的 56.5 美元（本书的核对）。
% 据: ps_link, ps_using

\fig{oil_price}{原油的年均价，1960—2025 年，每桶美元，对数刻度。名义价格与按 2010 年美元计的实际价格（用制造品出口单位价值指数折算）。阴影是 1973—1974 年与 1979—1980 年。数据：世界银行商品价格数据（Pink Sheet）年度工作簿，布伦特、迪拜与西得克萨斯三种原油等权平均。}{fig:oil_price}

第二次冲击来自伊朗。1978 年 12 月，OPEC 在阿布扎比开会，决定 1979 年把标准油价从每桶 12.70 美元平均提高 10\%，分几步走\cite{boj1986b}。Hamilton 说，1978 年秋天罢工蔓延到伊朗的石油业，到 1979 年 1 月伊朗的产量每天减少 480 万桶，约占当时世界产量的 7\%\cite{hamilton2011}。沙特等国增产，补上了其中约三分之一\cite{hamilton2011}。1979 年 6 月，OPEC 在日内瓦把标准油价从每桶 14.55 美元提到 18 美元\cite{boj1986b}。
% 据: bj_7812, hm_iran, hm_third, bj_june, bj_geneva

到 1981 年底，标准油价是每桶 34 美元，约为 1978 年底的 2.7 倍\cite{boj1986b}。BIS 1980 年说，这一次油价的实际涨幅与 1973—1974 年是同一个量级\cite{bis1980}。按世界银行的年均价，1980 年是每桶 36.9 美元\cite{pinksheet2026}。1986 年油价跌回 14.4 美元，这一段在第二次冲击之后（本书的核对）。
% 据: bj_34, bj_7812, b80_real, ps_link

\section{供给冲击}\label{sec:oil-model}

先回到\chref{ch:naturalrate}的式~\eqref{eq:naturalrate-curve}。需求过旺时，失业降到自然率以下，通胀上升，两者朝相反的方向走。油价冲击却让两者一起上升，这就是\textbf{滞胀}（stagflation），\chref{ch:phillips}讲过这个名字的来历。原因在成本。进口的原油变贵，在同样的失业水平上，企业的成本更高，定价也更高（本书的解读）。把这一点写进曲线，就多出一项
\begin{keyeq}
\pi_t=\pi^{e}_t-\alpha\,(u_t-u^{n})+s_t .\label{eq:oil-curve}
\end{keyeq}
$s_t$ 是\textbf{供给冲击}（supply shock），原油的实际价格上涨时为正。这是本书对\chref{ch:naturalrate}的曲线的推广（本书的推导）。原油是投入品，它变贵了，同样的产出要花更多的成本，企业会把一部分成本加到售价上。需求与预期都不变，物价也会多涨一截，这一截就是 $s_t$（本书的解读）。油价停在高处以后，成本不再上升，$s_t$ 回到零（本书的解读）。冲击也压低了实际工资能达到的水平。工人若要保住原来的实际工资，名义工资就得多涨，成本又被推高一轮（本书的解读）。

需求一侧由政策决定。设名义需求每年增长 $x_t$，它等于通胀加上实际增长。实际增长低于趋势 $\bar g$，失业就上升，仿照\chref{ch:phillips}里奥肯的经验关系，写成 $u_t-u_{t-1}=-\beta\,(x_t-\pi_t-\bar g)$（教学用的假设）。预期沿用\chref{ch:naturalrate}的适应性预期，$\pi^{e}_t=\pi_{t-1}$。

政策有两个极端。\textbf{不迁就}（non-accommodation）是名义需求的增长不变。冲击一来，通胀上升，实际增长就下降，失业上升，此后失业高于自然率，通胀逐年回落（本书的推导）。名义需求的盘子不变，物价涨得快了，能买的东西就少了，通胀的上升本身就起了紧缩的作用（本书的解读）。\textbf{迁就}（accommodation）是提高名义需求的增长，让失业留在自然率。这时 $\pi_t=\pi_{t-1}+s_t$，冲击过去以后，通胀也停在高处（本书的推导）。把式~\eqref{eq:oil-curve}逐年相加，得到
\begin{keyeq}
\sum_{t=1}^{T}(u_t-u^{n})=\frac{\pi_0-\pi_T+\sum_{t=1}^{T}s_t}{\alpha} .\label{eq:oil-choice}
\end{keyeq}
与式~\eqref{eq:naturalrate-sacrifice}相比，冲击给要压下去的通胀又添了一笔 $\sum s_t$（本书的推导）。要么多失业，要么让通胀永久升高，或者两者各担一部分，哪一种政策都免不了（本书的推导）。两个极端之间还有中间的路。每让通胀多留 1 个百分点，累计的失业就少 $1/\alpha$ 个百分点·年（本书的推导）。

这个式子记的是一笔账。供给冲击把物价水平抬高一截，适应性预期把这一截带进下一年的通胀。要让通胀回到原处，失业就得高于自然率，替这一截付账（本书的解读）。人们定预期的办法不同，账也会不同，后面误区框里 Gordon 的几种情形就差在这里（本书的解读）。

油价上涨还有一层损失，与曲线无关。BIS 1980 年说，油价上涨向进口国征了一笔税，减少它们的实际收入，因而压低支出\cite{bis1980}。进口量不变，这笔税约是
\begin{equation}
\frac{\Delta Y}{Y}\approx-\omega\,\frac{\Delta p_{o}}{p_{o}} ,\label{eq:oil-tot}
\end{equation}
$\omega$ 是净石油进口占国民生产总值的比例，$p_{o}$ 是油价（本书的推导）。税基是石油进口占收入的比例，税率是油价的涨幅，进口国越依赖石油，交的税越多（本书的解读）。按 BIS 的表，1973 年这个比例在联邦德国是 1.6\%\cite{bis1974}。欧洲的到岸价涨了约 144\%，多付的钱约合国民生产总值的 $1.6\times1.44\approx2.3\%$（本书的计算）。CIA 估计西欧与日本约 3\%\cite{frus47}。日本银行算出，1974 年日本进口原油多付了 124 亿美元，约合当时国民总支出的 3\%\cite{boj1986b}。
% 据: b80_tax, b74t_head, b74t_de, b74_cif, f277_loss, bj_transfer

式~\eqref{eq:oil-tot}与曲线说的是两件事。曲线说物价怎样变，这个式子说国民收入少了多少（本书的解读）。进口国为同样多的石油付出更多的出口品，这笔差额转给了产油国，货币政策要不回来（本书的解读）。BIS 1980 年说，私人部门要夺回被油价侵蚀的实际收入，当局要决定迁就多少、抵抗多少\cite{bis1980}。迁就会让工资与物价轮番上涨，抵抗会让失业增加，政策只能决定损失以哪种形式出现（本书的解读）。
% 据: b80_painful

\chref{ch:naturalrate}发现，失业在 7\% 以上的 1977、1980 与 1984 年，美国的通胀仍在上升。按世界银行的年均价，1980 年油价比上年涨了约 19\%，1977 年只涨了约 9\%，实际价格没有变，1984 年油价还跌了（本书的核对）。油价只说明得了 1980 年，另外两年的 $s_t$ 要另找来源（本书的解读）。
% 据: ps_link

\begin{compare}{需求冲击与供给冲击}{需求冲击}{供给冲击}
\cmprow{通胀与失业}{反向变动，需求过旺时通胀升、失业降（本书的推导）}{同向变动，$s_t>0$ 时通胀与失业一起上升（本书的推导）}
\cmprow{紧缩货币的效果}{通胀回落，失业回到自然率，两个目标一致（本书的推导）}{通胀回落，失业更高，两个目标冲突（本书的推导）}
\cmprow{实际收入}{冲击过去就恢复（本书的推导）}{付给产油国的部分收不回来，约合 GNP 的 2\% 到 3\%\cite{frus47,boj1986b}}
\cmprow{当时的看法}{CIA 1973 年 10 月认为就业起初不受影响\cite{frus47}}{CIA 1974 年 1 月担心产出下降与失业增加\cite{frus47}}
\cmpsum{两种冲击都推高通胀。只有在 $s_t=0$ 时，稳定通胀与稳定就业才是同一件事（本书的推导）。}
% 据: f223_jobs, f277_unless, f277_loss, bj_transfer
\end{compare}

\begin{example}{迁就与不迁就}\label{ex:oil-paths}
设 $\alpha=0.5$，$\beta=0.5$，$u^{n}=6\%$，$\bar g=3\%$，名义需求原来每年增长 7\%，通胀稳定在 4\%（教学用的假设）。第 1 年来了一次冲击，$s_1=4$，以后 $s_t=0$。（1）不迁就，名义需求仍增长 7\%，头两年的通胀与失业是多少？（2）完全迁就，通胀停在多少？名义需求要增长多少？（3）两种政策累计多出的失业各是多少？只让通胀最后停在 6\%，又是多少？
\solution
（1）第 1 年 $\pi_1=4-0.5\,(u_1-6)+4$，$u_1=6+0.5\,(\pi_1-4)$，解得 $\pi_1=7.2$，$u_1=7.6$。第 2 年同样解得 $\pi_2=5.92$，$u_2=8.56$。此后通胀继续回落，失业慢慢回到 6\%（本书的计算）。（2）失业守在 6\%，通胀从 4\% 升到 $4+4=8\%$ 并停在那里，名义需求要增长 $8+3=11\%$。（3）按式~\eqref{eq:oil-choice}，不迁就是 $(4-4+4)/0.5=8$ 个百分点·年，迁就是 0。通胀停在 6\% 时是 $(4-6+4)/0.5=4$ 个百分点·年（本书的计算）。
\end{example}

\fig{oil_paths}{\exref{ex:oil-paths}的路径（本书的计算）。(a) 通胀；(b) 失业率。蓝色实线是不迁就，名义需求每年增长 7\%；琥珀色虚线是完全迁就，失业守在自然率 6\%。第 1 年的供给冲击 $s_1=4$。参数是教学用的假设。}{fig:oil_paths}

\figref{fig:oil_paths}画出两条路。不迁就，第 1 年通胀升到 7.2\%，失业升到 7.6\%（本书的计算）。第 3 年失业最高，约 8.8\%（本书的计算）。此后通胀降到 4\% 以下，失业在 6\% 上下摆动着回来（本书的计算）。迁就，失业一直是 6\%，通胀从此是 8\%（本书的计算）。失业为什么会摆动？通胀降到 4\% 以下以后，名义需求仍增长 7\%，实际增长就高于趋势，失业又被压到自然率以下（本书的解读）。

\begin{pitfall}[油价上涨一定带来持久的通胀]
不一定。Gordon 1975 年分析了几种情形\cite{gordon1975}。他说，非农部门的预期若是适应性的，迁就的政策使物价水平永久升高，通胀率却不会\cite{gordon1975}。供给的减少若是永久的，比如牢不可破的石油卡特尔造成的，迁就才会使通胀率永久升高\cite{gordon1975}。本书的例题里，人们按上一年的通胀定预期，迁就一次冲击，通胀就停在高处，这比 Gordon 的情形更极端（本书的解读）。
% 据: gd_level, gd_cartel
\end{pitfall}

\section{迁就还是不迁就}\label{sec:oil-policy}

进口国面对两难。BIS 1980 年说，压缩需求来对付通胀，有加重石油冲击紧缩作用的危险。维持国内需求，又会助长通胀\cite{bis1980}。它说，迁就能减轻衰退，却有通胀永久升高的危险\cite{bis1980}。抵抗能压住通胀，却有经济下滑的危险\cite{bis1980}。
% 据: b80_choice, b80_painful

谁作的决定？两次冲击的选择不同。BIS 说，1973—1974 年主要工业国远没有一致把对付通胀放在首位，多数国家至少有一段时间选择了支持国内需求\cite{bis1980}。第二次冲击时，各国的货币政策明显偏紧，连最需要紧缩的美国也是这样，财政政策也在收紧\cite{bis1980}。美联储这几年的紧缩在\chref{ch:credibility}。CIA 1974 年 1 月也预料，进口国宁可让国际收支恶化，也要维持增长与就业\cite{frus47}。
% 据: b80_unanimous, b80_restrictive, b80_us, b80_fiscal, f277_choose

BIS 1975 年的年报记下了第一次的做法。1974 年，荷兰与瑞典先给了一些财政刺激，加拿大与英国的预算也在这一年变得更有扩张性\cite{bis1975}。衰退加深以后，货币政策普遍放松，动作最大的是联邦德国与美国\cite{bis1975}。可是许多国家通胀与国际收支的问题还在，BIS 说它们因此很谨慎\cite{bis1975}。
% 据: b75_fiscal, b75_eased, b75_caution

\begin{compare}{迁就与不迁就}{迁就}{不迁就}
\cmprow{名义需求}{增长加快，跟上更高的物价（本书的推导）}{增长不变，通胀的上升起紧缩的作用（本书的推导）}
\cmprow{BIS 1980 年的说法}{能减轻衰退，却有通胀永久升高的危险\cite{bis1980}}{能压住通胀，却有经济下滑的危险\cite{bis1980}}
\cmprow{例题的结果}{失业不变，通胀从 4\% 升到 8\% 并停住（本书的计算）}{累计多出 8 个百分点·年的失业，通胀回到原处（本书的计算）}
\cmprow{谁这样做}{1973—1974 年多数主要工业国至少一度支持国内需求\cite{bis1980}}{1979—1980 年各国的货币政策明显偏紧\cite{bis1980}}
\cmpsum{哪一种都要不回付给产油国的收入，两者只是把损失分成通胀与失业的不同比例（本书的解读）。}
% 据: b80_painful, b80_unanimous, b80_restrictive
\end{compare}

当时的理论也在找答案。Gordon 1975 年说，1973—1974 年粮食与石油的减产压低了美国非农工人的实际收入，又抬高了通胀\cite{gordon1975}。他认为，这种通胀与需求过旺的通胀不同，政策的对策也可以不同\cite{gordon1975}。工资与物价完全僵硬时，他说最优的政策是提高名义收入，迁就更高的外部价格，也保住原来的内部产出\cite{gordon1975}。
% 据: gd_title, gd_unnatural, gd_garden, gd_rigid, gd_optimal

到了 1980 年，BIS 的看法很明确。它说，1970 年代想靠追求增长走出滞胀的陷阱，最后通胀与失业都更糟，增长与通胀之间没有长期的交替\cite{bis1980}。它也说，除了很少几个例外，1979 年的货币工资与实际工资都比 1970 年代前半期涨得慢\cite{bis1980}。工资与物价普遍螺旋上升的危险，比 1973—1974 年小\cite{bis1980}。工资不追着物价涨，供给冲击就不容易变成持久的通胀，这是式~\eqref{eq:oil-curve}里预期那一项的作用（本书的解读）。
% 据: b80_growth, b80_tradeoff, b80_wages, b80_spiral

\figref{fig:oil_g7}画出七国集团与瑞士 1970—1985 年的通胀与增长。1974 年消费物价涨得最多的是日本，23.2\%，其次是意大利的 19.2\%\cite{wdi47}。联邦德国只有 7.0\%，瑞士是 9.8\%\cite{wdi47}。英国 1975 年最高，是 24.2\%\cite{wdi47}。增长按 Penn World Table 的实际 GDP 计，1974 年日本下降 1.2\%，英国下降 2.5\%，美国下降 0.5\%（本书的计算）\cite{pwt1001}。
% 据: wdi_cpi, pwt_rgdpna

\fig{oil_g7}{七国集团与瑞士的消费物价涨幅与实际 GDP 增长率（\%），1970—1985 年。实线是消费物价的年涨幅（世界发展指标），虚线是 Penn World Table 10.01 按不变本国价格计的实际 GDP 的增长率（本书的计算）。点与数字标出各国通胀最高与增长最低的一年。阴影是 1973—1974 年与 1979—1980 年。}{fig:oil_g7}

第二次冲击后，日本的通胀低得多，1979 年是 3.7\%，1980 年是 7.8\%\cite{wdi47}。意大利与英国 1980 年仍有 21.1\% 与 18.0\%\cite{wdi47}。联邦德国与瑞士的通胀，最高分别是 1981 年的 6.3\% 与 6.5\%，美国 1980 年是 13.5\%\cite{wdi47}。两国从 1975 年起公布货币增长的目标，这是\chref{ch:monetarism}的题目。衰退却来得晚些。按 Penn World Table，美国 1982 年下降 1.8\%，加拿大下降 3.2\%（本书的计算）\cite{pwt1001}。BIS 没有给各国排名次，数字也只说明结果，说明不了每一国作了什么选择（本书的解读）。
% 据: wdi_cpi, pwt_rgdpna

瑞士下降得最急。按 Penn World Table，瑞士 1975 年的实际 GDP 下降 7.3\%（本书的计算）\cite{pwt1001}。\chref{ch:monetarism}引 Rich 说产出下降 7.5\%，那是石油冲击加上紧缩的货币政策的结果\cite{rich2003}。两个数量的是同一件事，Rich 没有说明他用的产出序列，Penn World Table 用的是按 2017 年不变本国价格计的实际 GDP（本书的核对）。瑞士的通胀 1975 年降到 6.7\%，1979 年是 3.6\%\cite{wdi47}。通胀低的国家也没有躲过衰退（本书的解读）。
% 据: pwt_rgdpna, wdi_cpi, rich_drop

日本的高速增长也在这时到了头。\chref{ch:eastasia}引 Krugman 的数，1973 年以前的十年里日本实际 GDP 每年增长 8.9\%，1973—1992 年人均 GDP 每年只增长 3\%。第一次石油危机正落在这个转折点上（本书的解读）。日本的应对也与多数国家不同。BIS 说多数主要工业国一度支持需求，日本银行却在 1973 年 12 月把公定利率一次提高 2 个百分点\cite{bis1980,boj1986b}。第二次，它在 1979 年 4 月就提高利率，比 6 月日内瓦的提价还早（本书的核对）\cite{boj1986b}。行史里的经过，见下面的案例框。
% 据: b80_unanimous, bj_1221, bj_0416, bj_june, bj_geneva

代价落在谁身上？首先是进口国的消费者，他们为同样的油多付了钱，这是 BIS 说的那笔税\cite{bis1980}。其次是失业的人。CIA 在 1974 年 1 月就预料到失业会增加\cite{frus47}。再次是不产油的发展中国家。美国驻伊朗大使在电报里说，发展中世界受这次提价的打击最重，印度的五年计划要被扔进废纸篓\cite{frus47}。
% 据: b80_tax, f277_unless, f271_india

\begin{case}{日本 · 1973—1980 · 两次石油危机}
日本银行在行史里记下了两次危机。1973 年，大藏大臣爱知揆一（1907—1973）猝死，11 月下旬内阁改组，福田赳夫（1905—1995）接任大藏大臣，政府首脑基本上倾向加强紧缩与优待储蓄\cite{boj1986b,loc47}。12 月 21 日，日本银行开临时政策委员会，决定把公定利率提高 2 个百分点，商业票据贴现率从年 7\% 提到 9\%，22 日实行\cite{boj1986b}。公定利率是央行贴现商业票据的利率（本书的解读）。日本银行说，贴现率到 9\%，是明治 19 年（1886 年）以后的第一次，一次提高 2 个百分点也是从来没有过的\cite{boj1986b}。
% 据: bj_fukuda, bj_aichi, loc_aichi, bj_1221, bj_meiji, bj_never, loc_fukuda

人心也乱了。12 月 14 日，爱知县的丰川信用金库发生挤兑\cite{boj1986b}。1974 年 1 月 12 日，福田在大阪对记者说，物价“简直是狂乱状态”，“狂乱物价”从此成了通用的说法\cite{boj1986b}。批发物价到 2 月比上年同月上涨 37.0\%，1974 年平均上涨 31.3\%\cite{boj1986b}。消费物价 1974 年平均上涨 24.5\%\cite{boj1986b}。实际国民总支出比上年减少 1.2\%，是 1947 年以来第一次负增长\cite{boj1986b}。
% 据: bj_toyokawa, bj_kyoran, bj_wpi, bj_wpiavg, bj_cpi, bj_gne

代价落在谁身上？日本银行说，石油提价恶化了贸易条件，实际购买力转移到了产油国等海外\cite{boj1986b}。1974 年的春季工资谈判，基本工资提高 32.9\%，主要企业的夏季奖金比上年多 47\%，企业轻易答应了\cite{boj1986b}。完全失业者原来是 60 万到 70 万人，1975 年起急增到 100 万人上下\cite{boj1986b}。1974 年 12 月 16 日，总裁任满退任，次日森永贞一郎（1910—1986）接任\cite{boj1986b,loc47}。
% 据: bj_transfer, bj_wage, bj_easy, bj_unemp, bj_governor, ndl_morinaga

第二次危机，日本银行先动了手。1978 年 4 月初，它定下方针，货币供给与贷款若有越过危险线的苗头，即使在不景气中也要及时转向紧缩，不重犯“昭和 46—47 年的失败”\cite{boj1986b}。1979 年 1 月 23 日，森永说今后不再加大放松的力度\cite{boj1986b}。行史说，日本银行内部对上一次的形势判断有很强的反省，认为当时“过于乐观”，政策迟了\cite{boj1986b}。4 月 16 日，它把公定利率从 3.5\% 提到 4.25\%\cite{boj1986b}。
% 据: bj_apr78, bj_failure, bj_morinaga, bj_amaki, bj_0416, bj_0417

结果不一样。上一次货币供给比上年增长 20\% 多，这一次比较温和\cite{boj1986b}。消费物价的涨幅，1974 年 2 月的高峰是 26.3\%，这一次的高峰是 1980 年 9 月的 8.9\%\cite{boj1986b}。实际国民总支出 1979 年增长 5.2\%，1980 年增长 4.8\%\cite{boj1986b}。日本银行给的原因包括，危机前的经济情况不同，货币政策动手早，工资涨幅没有大大偏离劳动生产率的趋势\cite{boj1986b}。按它的计算，石油多付的钱占名义国民总支出的比例，第一次是 2.0\%，第二次是 2.2\%，冲击的规模差不多\cite{boj1986b}。
% 据: bj_ms20, bj_cpi89, bj_t35, bj_initial, bj_early, bj_wages, bj_ratio
\end{case}

\section{石油美元流向哪里}\label{sec:oil-recycling}

产油国一时花不完这笔钱。CIA 1974 年 1 月估计，产油国这一年的石油收入约 950 亿美元，是上一年的 3.5 倍\cite{frus47}。这个估计假定油价全年停在当时的水平，世界石油出口与 1973 年差不多，约每天 3200 万桶\cite{frus47}。沙特与伊朗各多收约 140 亿美元，委内瑞拉约 80 亿美元，科威特与利比亚各约 50 亿美元\cite{frus47}。它说，沙特、科威特与其他海湾小国的收入，会远远超过它们花钱的能力\cite{frus47}。BIS 1975 年的数字相近，OPEC 原油出口的平均价格从 1973 年的每吨约 20 美元升到 1974 年的 77 美元，石油出口额从约 320 亿美元升到 1150 亿美元\cite{bis1975}。
% 据: f277_95, f277_assume, f277_receipts, f277_spend, b75_ton, b75_value

盈余与赤字是同一笔账的两面。BIS 说，OPEC 国家的经常账户盈余 1974 年约 700 亿美元，上一年不到 50 亿美元\cite{bis1975}。代价落在谁身上？不产油的发展中国家的经常账户赤字从约 80 亿美元升到 250 亿美元以上，BIS 说它们的负担重得多\cite{bis1975}。它们只能借钱，或者少进口（本书的解读）。
% 据: b75_surplus, b75_ldc

这笔盈余要有人借出去。BIS 说，1974 年产油国的储备增加 364 亿美元\cite{bis1975}。它们在美国的美元储备增加约 100 亿美元，在英国的英镑储备增加约 50 亿美元，其余的绝大部分放进了欧洲货币市场\cite{bis1975}。欧洲货币市场是存放在本国以外的美元等货币的市场（本书的解读）。回流的路是这样的，产油国把美元存进伦敦等地的银行，银行再把钱贷给有赤字的国家（本书的解读）。
% 据: b75_reserves, b75_placement

美国的银行也加入进来。1974 年 1 月下旬，美国取消了对资本流出的管制\cite{bis1975}。BIS 说，这一年美国银行对外国的贷款猛增 177 亿美元，在把资金从盈余地区转到赤字地区时出了大力\cite{bis1975}。回流让进口国借钱付油账，当年的紧缩轻了一些，债务却留了下来（本书的解读）。
% 据: b75_controls, b75_banks

第二次比第一次难。BIS 1980 年警告，通过国际银行业的回流可能比 1973 年提价以后更难\cite{bis1980}。它说，与 1974 年相比，银行不再是从干净的账本开始，已经背着沉重的第三世界风险\cite{bis1980}。1974 年以后，几类国家对国际金融市场的债务负担大大加重了\cite{bis1980}。按 BIS 的统计，不产油的发展中国家的中期公共与公共担保外债里，欠私人债权人的比例从 1973 年的三分之一升到 1978 年的约一半\cite{bis1980}。对 1974 年以来风险暴露增长很快的一些国家，银行放贷可能不如以前积极\cite{bis1980}。这些贷款怎样变成浮动利率银团贷款，又怎样走到 1982 年的墨西哥危机，在\chref{ch:latamdebt}。
% 据: b80_harder, b80_burdened, b80_debt, b80_private, b80_keen

\section{荷兰病}\label{sec:oil-dutch}

油田带来的钱也会伤害产油国自己的经济。Corden 与 Neary 1982 年的模型把一国的经济分成三个部门\cite{cordenneary1982}。能源部门正在繁荣，制造业不景气，两者的产品都在世界市场上按国际价格买卖。第三个是服务业，产品不能贸易，价格在国内决定\cite{cordenneary1982}。他们用不可贸易品对可贸易品的相对价格来定义实际汇率\cite{cordenneary1982}。他们举的例子有澳大利亚的矿产、荷兰的天然气，以及英国、挪威与一些 OPEC 成员的石油，受压的都是传统的制造业\cite{cordenneary1982}。
% 据: cn_three, cn_rer, cn_cases

下面用本书的记号。三个部门是能源 E、制造业 M 与服务业 S。制造品当作计价单位，服务的相对价格记作 $p_S$。$p_S$ 上升，就是这种意义上的实际汇率升值。它对应\chref{ch:external}的式~\eqref{eq:external-real}里 $q$ 的下降，本国的东西相对外国变贵了（本书的解读）。
% 据: cn_rer

繁荣通过两条路影响别的部门\cite{cordenneary1982}。一条是\textbf{资源转移效应}（resource movement effect），能源部门的劳动边际产品提高了，从制造业与服务业把劳动吸走\cite{cordenneary1982}。另一条是\textbf{支出效应}（spending effect），收入增加了，一部分花在服务上，服务涨价，劳动又从制造业流向服务业\cite{cordenneary1982}。劳动直接从制造业流向能源，他们叫作直接的去工业化\cite{cordenneary1982}。繁荣引起的实际升值又把制造业的产出进一步压低，他们叫作间接的去工业化\cite{cordenneary1982}。
% 据: cn_effects, cn_direct, cn_indirect

两条路对制造业是同一个方向，对服务业的方向相反。Corden 与 Neary 说，服务业的产出是升是降，全看两种效应哪个强\cite{cordenneary1982}。他们认为英国的情形主要是支出效应\cite{cordenneary1982}。按世界发展指标，英国燃料出口占商品出口的比例从 1970 年的 2.6\% 升到 1982 年的 20.3\%\cite{wdi47}。这些年英国制造业的比例，世界发展指标没有数。Hutchison 的分解里，英国制造业的波动来自能源价格冲击的部分是三国里最多的\cite{hutchison1990}。Corden 与 Neary 也提醒，英国制造业的比例本来就因为别的原因在下降，能源繁荣的影响叠加在这个趋势上\cite{cordenneary1982}。
% 据: cn_hinges, cn_britain, wdi_fuel, hu_t4, hu_t4_uk, cn_trend

\begin{example}{能源繁荣与三个部门}\label{ex:oil-dutch}
设三个部门的产出都是 $X_i=A_i\sqrt{L_i}$，劳动总量为 1，可以自由流动，收入的一半花在服务上（教学用的假设）。$A_M=A_S=1$，能源的价值生产率 $p_EA_E$ 从 0.5 升到 1。（1）繁荣前后，三个部门各用多少劳动？$p_S$ 升多少？（2）$p_S$ 先不变，只有资源转移效应，劳动怎样分配？（3）按\chref{ch:external}的实际汇率，$q$ 降多少？
\solution
工资相同，$L_i$ 与 $(p_iA_i)^{2}$ 成正比。服务的开支等于收入的一半，可以解出 $L_S=0.5$，$p_S^{2}$ 等于另两个部门的 $(p_iA_i)^{2}$ 之和（本书的推导）。（1）繁荣前，$p_S^{2}=0.25+1=1.25$，$L_E=0.25/2.5=0.1$，$L_M=1/2.5=0.4$。繁荣后，$p_S^{2}=1+1=2$，$L_E=L_M=0.25$。$p_S$ 从 1.118 升到 1.414，升了约 26.5\%（本书的计算）。（2）$p_S^{2}$ 仍是 1.25，三项之和是 $1.25+1+1=3.25$，$L_E=L_M=1/3.25\approx0.308$，$L_S=1.25/3.25\approx0.385$（本书的计算）。（3）设物价指数里服务的权重是一半，$q$ 与 $p_S^{-0.5}$ 成正比，$q$ 降 $1-1/\sqrt{1.265}\approx11.1\%$（本书的计算）。
\end{example}

\begin{table}[tbp]\centering\small
\caption{\exref{ex:oil-dutch}里的两种效应：三个部门的劳动、服务的相对价格与工资}\label{tab:oil-dutch}
\begin{tabularx}{\linewidth}{@{}lYYYYYY@{}}
\toprule
 & 能源 $L_E$ & 制造业 $L_M$ & 服务业 $L_S$ & $p_S$ & 工资 & 制造业产出\\
\midrule
繁荣前 & 0.100 & 0.400 & 0.500 & 1.118 & 0.791 & 0.632\\
只有资源转移效应 & 0.308 & 0.308 & 0.385 & 1.118 & 0.901 & 0.555\\
再加上支出效应 & 0.250 & 0.250 & 0.500 & 1.414 & 1.000 & 0.500\\
\bottomrule
\end{tabularx}
\datasource{本书的计算，参数是教学用的假设。工资按制造品计。服务业的劳动最后回到 0.5，是这组函数形式的特例，Corden 与 Neary 说服务业可升可降。}
\end{table}

\tabref{tab:oil-dutch}把两种效应分开。资源转移效应把制造业的劳动从 0.4 减到约 0.3，服务业也少了。支出效应使服务涨价，服务业又把劳动拉回去，制造业再减到四分之一（本书的计算）。制造业的产出共减少约 21\%，一半多来自资源转移效应（本书的计算）。
% 据: cn_effects, cn_hinges

表里的工资也在变。按制造品计，工资从约 0.8 升到 1，涨了约 26.5\%，与服务价格涨得一样多（本书的计算）。制造品的价格由世界市场定，工资涨了，制造业的利润就被挤压，这是它收缩的原因（本书的解读）。从外面看，这个国家的东西变贵了，出口制造品更难（本书的解读）。能源繁荣不一定是坏事，国民收入是增加的，变的是收入从哪些部门来（本书的解读）。Corden 与 Neary 自己也指出，他们假定繁荣带来的收入由直接得益的人花掉\cite{cordenneary1982}。实际上能源部门的租金很大一部分以税的形式交给政府，政府怎样花这笔钱，决定了支出效应的大小与方向\cite{cordenneary1982}。花多少、花在哪里、什么时候花，是政府的决定，尼日利亚就是这样（本书的解读）。
% 据: cn_effects, cn_taxes

\begin{case}{荷兰 · 1959—1982 · 格罗宁根的天然气}
病名来自荷兰。Hutchison 在 BIS 的工作论文里说，荷兰 1959 年在格罗宁根省发现了巨大的斯洛赫特伦天然气田，十年后英国与挪威在北海发现了石油\cite{hutchison1990}。世界银行的 Budina 等人说，荷兰病得名于荷兰在 1960 年代管理天然气财富的不良记录\cite{budina2007}。
% 据: hu_slochteren, bu_named

数据上看得到病征。按世界发展指标，荷兰的制造业增加值占 GDP 的比例从 1970 年的 21.6\% 降到 1982 年的 14.9\%，燃料出口占商品出口的比例从 10.7\% 升到 23.9\%\cite{wdi47}。挪威的制造业比例从 17.8\% 降到 11.9\%，燃料出口从 0.9\% 升到 52.3\%\cite{wdi47}。
% 据: wdi_manf, wdi_fuel

可这是不是病，Hutchison 的检验并不支持。他用向量自回归分解制造业增长的波动，10 个季度后，能源价格的冲击在荷兰解释 9\%，在挪威解释 5\%，在英国解释 23\%\cite{hutchison1990}。他的结论是，正式的统计分析只微弱地支持制造业衰退与能源繁荣有系统的关系，支持最多的是挪威\cite{hutchison1990}。在荷兰，他几乎没有发现天然气开发对制造业有系统而长期的净损害\cite{hutchison1990}。他说，找到的影响通常在资源开发的初期，时有时无，也是暂时的，各国差别很大\cite{hutchison1990}。
% 据: hu_t4, hu_t4_nl, hu_t4_no, hu_t4_uk, hu_weak, hu_nl, hu_episodic
\end{case}

\begin{case}{尼日利亚 · 1962—1983 · 石油繁荣}
尼日利亚的出口几年里变了样。按世界发展指标，燃料占商品出口的比例 1962 年是 10.3\%，1970 年是 58.1\%，1974 年是 93.0\%\cite{wdi47}。食品的比例从 1962 年的 64.5\% 降到 1974 年的 5.1\%，农业原料从 16.6\% 降到 1.0\%\cite{wdi47}。
% 据: wdi_fuel, wdi_food, wdi_agri

谁作的决定？Budina、Pang 与 van Wijnbergen 说，1973 与 1974 年的大笔盈余先存在了国外，此后政府的开支迅速增加，到 1976 年吃掉了全部的石油横财\cite{budina2007}。1977 年联邦与各州的资本开支是 1970 年的 6 倍\cite{budina2007}。他们说，许多项目仓促上马，大多以浪费与腐败告终\cite{budina2007}。
% 据: bu_lag, bu_sixfold, bu_waste

1980 年代初油价下跌，债务随之而来。1981—1983 年，公共与公共担保的外债从 43 亿美元增加到 112 亿美元，外汇储备从 100 亿美元降到 12.3 亿美元，几乎耗尽\cite{budina2007}。代价落在谁身上？他们说，开支的压力抬高了非贸易部门的相对价格与工资，资源从出口品（农业）流向服务业，传统出口行业衰落了\cite{budina2007}。
% 据: bu_debt, bu_shift
\end{case}

\section{资源诅咒的数据检验}\label{sec:oil-test}

资源多的国家长得慢吗？Sachs 与 Warner 1995 年用 1970—1989 年各国的数据作了回归\cite{sachswarner1995}。他们的资源依赖是期初初级产品出口占 GDP 的比例，系数是 $-6.920$，t 值是 $-3.28$\cite{sachswarner1995}。这个比例的标准差约是 13 个百分点，高出一个标准差，人均增长每年约低 0.9 个百分点\cite{sachswarner1995}。
% 据: sw_abstract, sw_period, sw_coef, sw_sd, sw_093, sw_sxp

本书用另一组数据重做一遍。资源依赖用世界发展指标的自然资源租金占 GDP 的比例，在每段时期的起点量（本书的计算）\cite{wdi47}。增长率是 Penn World Table 的人均实际 GDP 的对数增长率，再控制起点的人均收入（本书的计算）\cite{pwt1001}。起点穷的国家有追赶的余地，不控制它，资源的作用会与追赶混在一起（本书的解读）。租金是资源的产值减去开采成本（本书的解读）。Sachs 与 Warner 用出口，本书用租金，两者量的东西不一样。出口量的是一国卖出多少初级产品，租金量的是资源本身值多少，转口贸易多的国家出口的初级产品也多（本书的解读）。样本是两个数据集里都有数的经济体，海湾国家单独标出（本书的计算）。
% 据: wdi_rents, wdi_rents_def, pwt_rgdpna

结果取决于时期（\figref{fig:oil_curse}）。1970—1990 年有 114 个经济体，租金比例每高 1 个百分点，人均增长每年低约 0.1 个百分点，t 值约是 $-4.1$（本书的计算）。租金比例的标准差约 8.0 个百分点，高一个标准差，增长每年低约 0.9 个百分点，与 Sachs 和 Warner 的结果相近（本书的计算）。改用对异方差稳健的标准误，t 值约是 $-2.2$，仍然显著（本书的计算）。可这个回归只解释了各国增长差异的约 13\%，资源远不是增长的全部（本书的计算）。
% 据: wdi_rents, pwt_rgdpna

1990—2019 年有 159 个经济体，系数只有约 $-0.02$，高一个标准差只低约 0.3 个百分点（本书的计算）。去掉 6 个海湾国家，t 值降到约 $-1.6$，这个系数就不再显著（本书的计算）。1990 年海湾 6 国的租金都在 GDP 的 36\% 以上，此后人均增长每年在 $-1.5\%$ 到 1.7\% 之间，它们撑起了这一段的负系数（本书的计算）。1970—1990 年的样本里只有 4 个海湾国家，去掉它们，系数几乎不变（本书的计算）。
% 据: wdi_rents, pwt_rgdpna

散点图上，租金最高的几个国家多半长得慢。赞比亚 1970 年的租金约占 GDP 的 36.5\%，此后 20 年人均每年下降约 2.1\%（本书的计算）。沙特约 25.6\%，人均每年下降约 1.0\%（本书的计算）。可租金高的国家里也有长得快的，回归线说的只是平均（本书的解读）。
% 据: wdi_rents, pwt_rgdpna

\fig{oil_curse}{自然资源租金占 GDP 的比例（起点年）与此后人均实际 GDP 的年均增长率（\%）。(a) 1970—1990 年；(b) 1990—2019 年。三角是海湾国家，直线是控制起点人均收入以后的回归线，画在起点收入的平均值上（本书的计算）。数据：世界发展指标（NY.GDP.TOTL.RT.ZS）与 Penn World Table 10.01。}{fig:oil_curse}

也有反例。Sachs 与 Warner 说，1970—1992 年资源依赖最高的四分之一里，只有马来西亚与毛里求斯人均增长每年在 2\% 以上，博茨瓦纳也许也算，可惜缺数据\cite{sachswarner1995}。他们认为，这两国靠的是制造品出口\cite{sachswarner1995}。在本书的数据里，马来西亚 1970 年的租金约占 GDP 的 17\%，此后 20 年人均每年增长约 4.8\%（本书的计算）。
% 据: sw_exceptions, sw_manufacturing, wdi_rents, pwt_rgdpna

另两个常举的反例，在本书的数据里看不出来。博茨瓦纳 1970 年的租金只占 GDP 的约 3\%，此后 20 年人均每年增长约 9.0\%（本书的计算）。挪威 1970 年的租金不到 GDP 的 1\%，北海的石油那时刚发现不久（本书的解读）。到 1990 年，挪威的租金是 8.0\%，此后人均每年增长约 1.5\%（本书的计算）。在起点量租金，后来才繁荣的国家被算成资源少的国家（本书的解读）。
% 据: wdi_rents, pwt_rgdpna, hu_slochteren

回归说明不了因果。租金高的国家，也许政局不稳，也许制度差，回归分不开（本书的解读）。起点的租金与后来的油价有关，1970—1990 年包含了油价的暴涨与暴跌，1990—2019 年的油价大体在上涨，两段的系数不同，也许与此有关（本书的解读）。
% 据: ps_link

\section{争论}\label{sec:oil-debate}

\subsubsection{石油是不是滞胀的主因}

一方认为货币才是主因。Barsky 与 Kilian 2001 年说，油价大涨在 1970 年代滞胀的因果机制里，远没有通常以为的那么重要\cite{barskykilian2001}。他们说，理论上推不出石油供给冲击一定带来滞胀，经验上也没有稳健的证据\cite{barskykilian2001}。货币的扩张与收缩，即使没有供给冲击，也能产生大小相当的滞胀\cite{barskykilian2001}。\chref{ch:naturalrate}讲过，弗里德曼 1976 年也说光凭石油危机解释不了通胀与失业一起上升。
% 据: bk_abstract

另一方看重油价的记录。Hamilton 2011 年说，历次供给削减的时机、大小与性质都与地缘政治事件密切相关\cite{hamilton2011}。他还说，战后美国的 11 次衰退，除 1960 年的一次以外，都伴随着油价的上涨\cite{hamilton2011}。他也说，结论只是石油冲击至少在一部分衰退里是一个因素\cite{hamilton2011}。Bernanke、Gertler 与 Watson 1997 年认为，石油冲击本身不是那几次衰退的主要原因\cite{bgw1997}。货币政策对油价的内生反应，能解释油价冲击对实体经济的紧缩作用的很大一部分，有些情形里几乎是全部\cite{bgw1997}。他们说，那几次衰退里，货币政策与其他既非货币也非石油的扰动都起了重要作用\cite{bgw1997}。
% 据: hm_geo, hm_eleven, hm_some, bgw_perse, bgw_policy, bgw_both

2000 年代的经验又添了一层。按世界银行的实际价格，2008 年的油价是每桶 94.3 美元，高过 1980 年的 56.5 美元（本书的核对）。那几年却没有 1970 年代那样的滞胀。Blanchard 与 Galí 把对数油价累计上涨 50\% 以上、持续四个季度以上的情形算作大的石油冲击，得到 1973、1979、1999 与 2002 年开始的四次\cite{blanchardgali2007}。他们说，油价只解释了 1970 年代波动的一部分\cite{blanchardgali2007}。油价冲击的动态影响随时间大大减小了，冲击本身却没有变小\cite{blanchardgali2007}。他们给出 2000 年代不同的四个原因，没有同时来的其他不利冲击，石油在生产中的比重小了，劳动市场更灵活，货币政策改进了，并认为四个都起了重要作用\cite{blanchardgali2007}。至今没有定论。
% 据: ps_link, bg_define, bg_episodes, bg_part, bg_smaller, bg_four, bg_all

\subsubsection{资源诅咒是真是假}

Sachs 与 Warner 发现资源依赖高的国家增长慢，本书在 1970—1990 年的数据里也看到了，1990—2019 年就弱得多（\secref{sec:oil-test}）。Mehlum、Moene 与 Torvik 2002 年认为，关键在制度\cite{mmt2002}。他们说，制度有利于攫取者时，资源多会压低收入，有利于生产者时，资源多会提高收入\cite{mmt2002}。
% 据: mmt_grabber, sw_abstract

他们先承认平均的差距。资源多的尼日利亚、赞比亚、塞拉利昂、安哥拉、沙特与委内瑞拉都长得慢，资源少的韩国、台湾、香港与新加坡长得快\cite{mmt2002}。可资源多又长得好的也不少，他们举了博茨瓦纳、加拿大、澳大利亚与挪威\cite{mmt2002}。他们说，平均的差距还不足以让人直接断定有资源诅咒\cite{mmt2002}。
% 据: mmt_losers, mmt_winners, mmt_jump

荷兰病是不是那条传导的路，也有疑问。Budina 等人认为，拖累尼日利亚非石油经济的，是开支的剧烈波动，不是荷兰病效应\cite{budina2007}。他们说，在 90 个国家里，尼日利亚的贸易条件波动排第三，波动像是对贸易品生产投资的一种税\cite{budina2007}。Hutchison 在荷兰、挪威与英国只找到微弱的证据\cite{hutchison1990}。至今没有定论。
% 据: bu_volatility, bu_third, bu_tax, hu_weak

\begin{deeper}[史料说明]
FRUS 是美国国务院史官办公室编的档案集，本章读了 1969—1976 年第 36 卷的第 209、216、220、223、271、273、277 号文件与人名表，都是网页版。参加会议与发电报的官员，本章按人名表的职务称呼。基辛格的回忆录只按编者的转引。CIA 1973 年 10 月的备忘录与 1974 年 1 月的文件都只收了节录，没有收进 FRUS 的部分本书没有读。

几个价格在史料里写法不一。10 月 16 日的公示价格，日本银行记作 5.12 美元，BIS 1974 年的年报记作 5.11 美元。12 月 23 日的价格，日本银行与 BIS 都记作 11.65 美元，FRUS 的编者注记作 11.60 美元，正文用日本银行的数字。日本银行算出第二次冲击使油价涨了约 2.7 倍，指的是 34 美元与 12.70 美元之比，即涨到原来的 2.7 倍（本书的核对）。

BIS 的年报读的是 1974、1975 与 1980 年的扫描本及识别文字层，识别文字认错的地方照录，在快照的说明里一一列出。Hutchison 的工作论文没有文字层，日本银行的《日本银行百年史》第 6 卷也没有，都按渲染的页面图像读出。Barsky 与 Kilian 和 Sachs 与 Warner 的文字层错乱，引文按页面图像读出。Sachs 与 Warner 的起点年，摘要与正文一处说是 1971 年，另一处说是 1970 年，正文用 1970 年。

Corden 与 Neary 读的是 IIASA 的会议论文本，封面要求引用前征得作者同意，本章只转述，不引原文。它后来发表在《经济学杂志》上，本书没有取得那一版。Rich 说瑞士 1975 年的产出下降 7.5\%，Penn World Table 是 7.3\%。美国 1974 年的通胀，世界发展指标是 11.1\%，\chref{ch:naturalrate}按总统经济报告写作 11.0\%，四舍五入以前是 11.05\%（本书的核对）。日本银行的行史只写了 1974 年底退任的总裁的姓，本书没有找到能确认他名字与生卒年的规范记录，正文不写。

世界银行的年度工作簿印的更新日期是 2026 年 9 月 2 日，网页上的链接却写作 2026 年 11 月。工作簿的封面还列出两张表里没有的工作表。数值与 1974 年等年份的通行数字相合，本书照用（本书的核对）。世界发展指标的自然资源租金是石油、天然气、煤、矿产与森林租金之和。
% 据: fp_list, f271_387, b74_posted, b74_1165, f271_1160, bj_34, bj_7812, sw_abstract, sw_sxp, cn_cover, rich_drop, wdi_cpi, bj_governor, ps_link, wdi_rents_def
\end{deeper}

\section{今天还剩什么}\label{sec:oil-legacy}

供给冲击成了宏观经济学的常用词。实际经济周期理论把衰退看作对技术冲击的反应，在\chref{ch:rbc}。三方程模型里的供给冲击在\chref{ch:consensus}。2000 年代产油国的巨额盈余在\chref{ch:moderation}，2021—2023 年的能源冲击在\chref{ch:inflation21}，央行该不该看穿供给冲击，在\chref{ch:fragmentation}。

产油国的难题也还在。\chref{ch:greekroman}讲过雅典的劳里昂银矿，横财是分掉还是投资，这个问题今天的产油国还在回答（本书的解读）。\chref{ch:silver}讲过西班牙王室靠忽多忽少的白银过日子，尼日利亚的开支比石油收入波动还大，是同一类难处（本书的解读）\cite{budina2007}。油价也没有安定下来。按世界银行的实际价格，1998 年每桶只有 15.9 美元，2012 年是 95.3 美元，2020 年又跌到 42.1 美元（本书的核对）\cite{pinksheet2026}。
% 据: bu_more, ps_link

两次石油危机也改变了央行。日本银行第二次先动手，结果好得多，它自己把这归于记取了 1971—1972 年的教训（本书的概括）\cite{boj1986b}。BIS 1980 年也说，靠追求增长走不出滞胀（本书的概括）\cite{bis1980}。紧缩靠的是让人们相信通胀会降下来，这是\chref{ch:credibility}的题目。
% 据: bj_failure, bj_initial, b80_growth

\begin{takeaways}
\begin{itemize}
\item 1973 年 10 月，海湾产油国把公示价格提高 70\%，OAPEC 宣布每月减产 5\%。12 月 23 日的德黑兰会议又把价格提到 11.65 美元，沙特反对，伊朗国王带头。
\item 供给冲击使通胀与失业一起上升。式~\eqref{eq:oil-choice}说，冲击要么换成更多的失业，要么换成永久更高的通胀。付给产油国的实际收入，哪种政策都要不回来。
\item BIS 说，1973—1974 年多数工业国至少一度支持需求，第二次冲击时货币政策明显偏紧。日本第二次先紧缩，通胀高峰从 26.3\% 降到 8.9\%。
\item 1974 年产油国的经常账户盈余约 700 亿美元，大部分经欧洲货币市场与银行借了出去。BIS 1980 年警告，银行已经背着沉重的第三世界风险。
\item 荷兰病模型里，能源繁荣经资源转移效应与支出效应使制造业收缩，服务业的价格上升。资源诅咒在 1970—1990 年的数据里明显，在 1990—2019 年弱得多，制度与波动的作用至今有争论。
\end{itemize}
\end{takeaways}

\begin{selfcheck}
\begin{enumerate}
\item 按\exref{ex:oil-paths}的参数，若冲击是 $s_1=2$，不迁就时累计多出多少失业？完全迁就时通胀停在多少？
\item 净石油进口占 GNP 的 2\%，油价涨 150\%，进口量不变，实际收入约损失多少？
\item 按\exref{ex:oil-dutch}，若 $p_EA_E$ 从 0.5 升到 1.5，繁荣后 $L_M$ 是多少？
\item Sachs 与 Warner 的系数是 $-6.920$，资源依赖高出 0.05，人均增长每年约低多少？
\item 1973 年的公示价格从 3.01 美元提到 5.12 美元，涨了多少？
\end{enumerate}
\end{selfcheck}

\begin{answers}
\item $(4-4+2)/0.5=4$ 个百分点·年，通胀停在 $4+2=6\%$。
\item $2\times1.5=3\%$ 的 GNP。
\item $p_S^{2}=2.25+1=3.25$，$L_M=1/6.5\approx0.154$。
\item $6.920\times0.05\approx0.35$ 个百分点。
\item $5.12/3.01-1\approx70\%$。
\end{answers}

回到 1973 年 10 月 18 日的例会。基辛格问，油价涨到哪里才该抵抗？进口国的政府很快面对同样的问题，只是要抵抗的是通胀，代价是失业（本书的解读）。迁就的后果，取决于人们怎样形成对通胀的预期，Gordon 的几种情形就差在这里\cite{gordon1975}。人们若看穿了政策，用老办法估出来的关系也会跟着变（本书的解读）。人们会预期政策时，政策还有效吗？用历史数据估计的模型，为什么会在政策改变后失效？
% 据: f220_twenty, gd_level, gd_cartel
